\documentclass[10pt]{article}
\usepackage[letterpaper,margin=0.85in]{geometry}
\usepackage{booktabs}
\usepackage{microtype}
\usepackage{parskip}
\usepackage{enumitem}
\usepackage{hyperref}
\usepackage{xurl}
\hypersetup{colorlinks=true,linkcolor=black,urlcolor=blue,citecolor=black,pdftitle={Global and Stage Specific Concurrency Control for Agent Serving}}
\setlist{nosep}
\setcounter{secnumdepth}{0}

\newcommand{\makeproposalheader}{%
  \begin{center}
    \vspace*{-0.35in}
    {\LARGE\bfseries
      Global and Stage Specific\par
      \vspace{0.15em}
      Concurrency Control for Agent Serving\par}
    \vspace{0.8em}
    {\large\scshape Course Project Proposal\par}
    \vspace{1.15em}
    {\large
      \begin{tabular}{c@{\hspace{4.5em}}c}
        Lingxiao Yang & Chenrui Gao
      \end{tabular}\par}
    \vspace{0.8em}
    {CS 395T: Hardware--Software Co-design for Cloud and Datacenters\par}
    \vspace{0.15em}
    {Instructor: Jovan Stojkovic\par}
    \vspace{0.8em}
    {Fall 2026 \textperiodcentered{} Revised October 9, 2026\par}
  \end{center}
  \vspace{0.35em}
}


\usepackage{longtable,array,calc}
\usepackage{amsmath}
\providecommand{\tightlist}{\setlength{\itemsep}{0pt}\setlength{\parskip}{0pt}}
\newcounter{none}

\begin{document}
\makeproposalheader
\hypertarget{motivation-and-research-question}{%
\subsection{Motivation and research
question}\label{motivation-and-research-question}}

Tool-using agents repeatedly move between model inference and tool
execution. An LLM output selects the next tool and its arguments; the
tool result determines the next prompt, context length, and sometimes
the remaining execution path. Scheduling one stage therefore changes the
arrivals seen by the other. Separate concurrency limits provide
additional control, but do not remove these dependencies.

We ask: \textbf{on a single CPU/GPU server, when is global request
concurrency control sufficient, and when do separately adjustable stage
limits improve end-to-end performance? Can an online policy follow
changes in workload and resource availability?} CPU interference is the
main diagnostic focus. GPU utilization, inference queues, correctness,
and request latency are needed to evaluate its consequences; maximizing
CPU utilization alone is not the objective.

The outcome will be an empirical map of operating conditions and an
open-source serving testbed, together with a comparison of static and
adaptive policies. A finding that global control is sufficient in a
measured regime is useful evidence, not a failed attempt to demonstrate
a predetermined speedup.

\hypertarget{revised-hypotheses}{%
\subsection{Revised hypotheses}\label{revised-hypotheses}}

\textbf{H1: Resource balance changes the useful control policy.} At
fixed hardware, changing the available tool CPU budget or tool
parallelism can change which stage constrains throughput. Stage-specific
limits may help when they reduce contention without creating greater
admission delay or inference starvation. These effects must be evaluated
jointly, not subtracted as independently measured gains.

\textbf{H2: There is no assumed universal tool-time threshold.} Policy
benefit depends on CPU capacity, tool type/threading, model service
behavior, offered load, and trajectory structure. We will report the
region where stage-specific control helps, ties, or hurts. A tool-time
fraction is a descriptive measurement for a stated configuration, not an
invariant workload property or a sufficient predictor.

\textbf{H3: Online adaptation may help under changing demand.} With the
same observations and tuning budget, adaptive stage-specific control may
outperform adaptive global control and tuned static policies during
workload-mix or CPU-budget changes. We will measure recovery time,
oscillation, and overhead as well as steady-state performance.

\hypertarget{scope-and-feasibility}{%
\subsection{Scope and feasibility}\label{scope-and-feasibility}}

The core study uses one server, one GPU, one model, and local tools.
Remote CPU services, multi-node placement, and different GPU generations
are outside the required evaluation. A second GPU platform is optional
replication; reducing a CPU budget does not emulate a different GPU
architecture.

Cheap tools remain a control case. For an isolated serial trajectory
with unchanged service times, a small tool-time fraction bounds the
direct benefit of speeding up that stage. It does not bound
admission-control gains under queueing, batching, or interference. We
will first measure live stage demand and static limit sweeps. If
policies are indistinguishable, report the tested boundary with
uncertainty and use a lightweight adaptive implementation to quantify
overhead rather than expanding into a complex scheduler.

\hypertarget{single-server-setup-and-workloads}{%
\section{Single Server Setup and
Workloads}\label{single-server-setup-and-workloads}}

\hypertarget{execution-stack}{%
\subsection{Execution stack}\label{execution-stack}}

One Linux server hosts a load generator, a LangGraph ReAct-style
data-analysis agent, one shared \textbf{Qwen3-8B/vLLM} instance, a CPU
tool-executor service, and a trace collector. Model and tool calls use
localhost HTTP. Plan for about 16 CPU cores, 64 GB RAM, and one GPU with
24-48 GB VRAM; actual access and completed implementation are not
established by this proposal.

A shared executor uses a bounded process pool and FIFO tool queue.
Inference submissions pass through a separate admission queue before
vLLM. Resources for vLLM host threads, orchestration, and monitoring are
reserved and recorded separately from the tool CPU budget. The same
allocation is used for every policy comparison within a condition.

Pin model/tokenizer revisions, sampling configuration, context/token
limits, prompts, software, and fixtures. Validate service health,
tool-call parsing, gold outputs, trace IDs, and timeouts. Deploy with
Docker Compose or a cluster-compatible environment inside the assigned
allocation. If a smaller model is required, label it as a separate
configuration: model substitution can change decisions and token counts,
not just inference speed.

\hypertarget{verifiable-task-families}{%
\subsection{Verifiable task families}\label{verifiable-task-families}}

\begin{longtable}[]{@{}
  >{\raggedright\arraybackslash}p{(\columnwidth - 4\tabcolsep) * \real{0.3333}}
  >{\raggedright\arraybackslash}p{(\columnwidth - 4\tabcolsep) * \real{0.3333}}
  >{\raggedright\arraybackslash}p{(\columnwidth - 4\tabcolsep) * \real{0.3333}}@{}}
\toprule\noalign{}
\begin{minipage}[b]{\linewidth}\raggedright
Family
\end{minipage} & \begin{minipage}[b]{\linewidth}\raggedright
Tool computation and validation
\end{minipage} & \begin{minipage}[b]{\linewidth}\raggedright
Role
\end{minipage} \\
\midrule\noalign{}
\endhead
\bottomrule\noalign{}
\endlastfoot
Lightweight control & Indexed SQLite lookups and small aggregations;
reference SQL answers. & Identify inference-dominated cases and
controller overhead. \\
Data analysis & File parsing, joins, grouped aggregation, and numerical
statistics in a Python worker; reference numeric results and hashes. &
Candidate CPU/memory-intensive work; classify using measured
profiles. \\
\end{longtable}

Use 10 calibration and 20 held-out tasks per family. Choose fixture
sizes from declared task/data scales that fit the machine, retaining
normal indexes and data types. Freeze fixtures before policy tuning and
publish all sizes; do not select them to obtain a desired tool-time
fraction or imitate an unavailable production distribution. Preserve
errors, retries, and failed answers.

\hypertarget{controlled-resource-dimensions}{%
\subsection{Controlled resource
dimensions}\label{controlled-resource-dimensions}}

\begin{longtable}[]{@{}
  >{\raggedright\arraybackslash}p{(\columnwidth - 4\tabcolsep) * \real{0.3333}}
  >{\raggedright\arraybackslash}p{(\columnwidth - 4\tabcolsep) * \real{0.3333}}
  >{\raggedright\arraybackslash}p{(\columnwidth - 4\tabcolsep) * \real{0.3333}}@{}}
\toprule\noalign{}
\begin{minipage}[b]{\linewidth}\raggedright
Dimension
\end{minipage} & \begin{minipage}[b]{\linewidth}\raggedright
Core comparison
\end{minipage} & \begin{minipage}[b]{\linewidth}\raggedright
Interpretation
\end{minipage} \\
\midrule\noalign{}
\endhead
\bottomrule\noalign{}
\endlastfoot
Tool CPU availability & For example 2, 4, and 8 dedicated cores on the
same host & Changes effective CPU/GPU balance with model and GPU
fixed. \\
Within-tool parallelism & One versus a small fixed thread count for
tools that support it & Measures interaction between concurrent calls K
and threads per call P. \\
Workload and arrivals & Light/heavy mixtures, steady arrivals, one burst
and one mix shift & Tests policy robustness without changing the
model. \\
\end{longtable}

Thread settings remain fixed during each policy comparison. A two-thread
tool may not achieve twice the speed; measure scaling and CPU-seconds
rather than infer demand from K alone. Start with serial tool
dependencies per request. Parallel fan-out within one agent is optional
and must preserve dependency correctness. If affinity/cgroups are
unavailable, report the actual isolation method and omit unsupported
CPU-budget claims.

CPU and GPU capacities are profiled separately to explain the operating
point, then validated jointly in live serving. GPU request latency is
not the reciprocal of GPU request throughput because inference batches
overlap. Record input/output tokens, action counts, bytes returned, and
tool results so changes in execution path remain visible.

\hypertarget{policies-and-continuous-monitoring}{%
\section{Policies and Continuous
Monitoring}\label{policies-and-continuous-monitoring}}

\hypertarget{admission-controls}{%
\subsection{Admission controls}\label{admission-controls}}

\begin{longtable}[]{@{}
  >{\raggedright\arraybackslash}p{(\columnwidth - 2\tabcolsep) * \real{0.5000}}
  >{\raggedright\arraybackslash}p{(\columnwidth - 2\tabcolsep) * \real{0.5000}}@{}}
\toprule\noalign{}
\begin{minipage}[b]{\linewidth}\raggedright
Control
\end{minipage} & \begin{minipage}[b]{\linewidth}\raggedright
Definition
\end{minipage} \\
\midrule\noalign{}
\endhead
\bottomrule\noalign{}
\endlastfoot
N & Maximum active end-to-end agent requests, including requests waiting
between stages. \\
K & Maximum executing tool calls; waiting calls occupy no worker. \\
M & Maximum model calls submitted to vLLM and not yet complete,
including its internal waiting calls. \\
\end{longtable}

M is an external admission limit, not a GPU batch-size setting. A
request holds a stage slot only for that call and releases it before
entering another stage. Bound all queues; count admission waits and
rejected arrivals. Decreasing a limit does not preempt work already
running. Global policies tune N with nonbinding internal gates; stage
policies hold a common outer safety cap and adjust K and M. Internal
gates still obey fixed memory-safety limits.

Compare a tuned static global policy N, a tuned static stage policy (K,
M), an adaptive global policy, and an adaptive stage policy. Default
serving is a secondary reference. Add a K-only ablation at selected
conditions to determine whether an inference-side gate adds value. This
distinguishes the benefit of adaptation from the benefit of more control
points.

\hypertarget{online-controller}{%
\subsection{Online controller}\label{online-controller}}

Both adaptive policies use the same lightweight feedback algorithm,
initialization procedure, observation budget, and end-to-end objective:
improve SLO goodput subject to fixed per-family latency and correctness
requirements. The global policy can change N; the stage policy can
change K or M. Offline calibration establishes safe ranges and initial
values, while runtime measurements drive subsequent actions.

\begin{enumerate}
\def\labelenumi{\arabic{enumi}.}
\tightlist
\item
  \textbf{Observe continuously.} Update queue depth, occupancy,
  arrivals, and completions on each event; sample process CPU and
  GPU/server metrics periodically. Maintain rolling tool-type service
  times, queue waits, and completion/goodput estimates.
\item
  \textbf{Choose a bounded probe.} When the window has enough
  observations, change one limit by one step. Pressure and starvation
  signals guide which direction to test. Low GPU utilization alone never
  triggers an increase.
\item
  \textbf{Evaluate the coupled effect.} Wait for enough affected work to
  complete before accepting another change. Retain a change only if the
  end-to-end objective improves without violating guards; otherwise
  revert. Use hysteresis, minimum dwell time, and hard bounds to limit
  oscillation.
\end{enumerate}

This is one controller with separately adjustable stage limits, not two
independent controllers optimizing local utilization. Fine-grained event
monitoring does not require changing limits on every event. Decision
windows scale with measured stage durations; sparse samples defer
updates. Thresholds and window rules are fixed before held-out testing.
Online control uses deadline-compliant completions as a goodput proxy;
gold-answer correctness is checked separately by the benchmark and must
not regress.

\hypertarget{profiling-and-attribution}{%
\subsection{Profiling and attribution}\label{profiling-and-attribution}}

OpenTelemetry links arrival, inference admission/submission/completion,
tool enqueue/start/end, and response by request ID. py-spy and
process/cgroup counters measure CPU hotspots/demand; perf/scheduler
traces investigate contention where permitted. vLLM queue/token and GPU
metrics provide inference context. Log policy actions and their effects.

Avoid double-counting overlapping spans or assigning aggregate CPU usage
to one concurrent request. Separate tool computation, framework
overhead, and waiting; validate accounting and measure tracing overhead.
GPU idleness alone does not establish tool-induced starvation: require
matched interventions and live confirmation. Report counter visibility
limits.

\hypertarget{evaluation-and-open-source-artifact}{%
\section{Evaluation and Open Source
Artifact}\label{evaluation-and-open-source-artifact}}

\hypertarget{experiments-and-success-criteria}{%
\subsection{Experiments and success
criteria}\label{experiments-and-success-criteria}}

Start with a small static sweep over feasible N, K, and M values, rather
than a full Cartesian search. Profile the CPU budgets and tool-thread
settings on calibration tasks, then select representative
inference-limited, tool-limited, and balanced conditions if these
regimes actually occur. Use the same calibration allowance for all
policies and freeze selections before held-out evaluation.

At selected operating points, compare all four policies under steady
load and a workload-mix transition. Add one controlled CPU-budget
reduction/recovery on the same host if allocation control is available.
Keep the disturbance, absolute arrival rates, tasks, and total available
resources identical across policies; controller actions never change the
resource budget. Include an inference-heavy control even if it shows no
benefit.

Measure \textbf{SLO goodput, successful throughput, p50/p95 end-to-end
latency, and per-family correctness}. Set per-family SLOs before testing
and count time from scheduled arrival, including every admission queue.
Report rejected, timed-out, unfinished, and incorrect tasks. Also
measure tool CPU-seconds, service/queue time, inference queueing,
tokens, memory, and GPU activity. For transitions, report settling time,
cumulative SLO misses, and limit oscillations.

Use three independent runs at selected conditions, randomized
configuration order, and documented warmup/cache state. Target 500 or
more completions per run for p95-led claims when compute permits, but
sample count alone does not establish precision. Report uncertainty with
task/run-aware resampling and increase samples only when needed to
distinguish a material effect. If intervals remain wide, narrow the
claim. Compare tracing and controller overhead separately.

The main figures are policy-benefit maps across CPU budget and tool
parallelism, latency/goodput curves, and aligned timelines of resource
changes, queues, and policy actions. Report effect sizes and uncertainty
in place of one universal phi threshold. Hold workload/model fixed for
hardware-allocation comparisons; report path/token changes in live
policy comparisons.

\hypertarget{replay-and-simulation}{%
\subsection{Replay and simulation}\label{replay-and-simulation}}

Recorded model actions plus real tool execution support deterministic,
paired policy comparisons; verify actions and tool-result hashes and
fail visibly on divergence. Paced replay fixes model delays and
therefore removes part of the very feedback under study. It is useful
for debugging and isolating mechanisms, but cannot validate online
CPU-GPU interaction, inference batching, or overall live gains. Core
conclusions require the real local model. A calibrated trace-driven
queue simulation is optional sensitivity analysis, not an additional
required system.

\hypertarget{reusable-release}{%
\subsection{Reusable release}\label{reusable-release}}

\begin{longtable}[]{@{}
  >{\raggedright\arraybackslash}p{(\columnwidth - 2\tabcolsep) * \real{0.5000}}
  >{\raggedright\arraybackslash}p{(\columnwidth - 2\tabcolsep) * \real{0.5000}}@{}}
\toprule\noalign{}
\begin{minipage}[b]{\linewidth}\raggedright
Module
\end{minipage} & \begin{minipage}[b]{\linewidth}\raggedright
Released functionality
\end{minipage} \\
\midrule\noalign{}
\endhead
\bottomrule\noalign{}
\endlastfoot
workloads/ and configs/ & Fixtures, reference checks, resource/thread
settings, and arrival/mix traces. \\
tracing/ and replay/ & Stage schema, record/replay adapters, dependency
checks, and CPU-only demo. \\
policies/ & Shared admission hooks, global/stage baselines, online
controller, and event logs. \\
experiments/ and analysis/ & Launch scripts, measured example outputs,
confidence intervals, and figure generation. \\
\end{longtable}

Runs emit checked answers, trajectories, spans, request metrics, N/K/M
events, and hardware/software manifests. Provide a clean-checkout guide
and commands for replay, live serving, and plots; CI covers accounting,
slot release, limit changes, and replay equivalence. Release original
code under Apache-2.0 subject to dependency compatibility, retain
upstream attribution, and publish sanitized fixtures/traces and model
acquisition instructions. These are planned deliverables, not completed
results.

\hypertarget{schedule-and-research-context}{%
\section{Schedule and Research
Context}\label{schedule-and-research-context}}

\hypertarget{remaining-semester-plan}{%
\subsection{Remaining semester plan}\label{remaining-semester-plan}}

Planned milestones below use Austin local time; record actual progress
and access in the midterm report.

\begin{longtable}[]{@{}
  >{\raggedright\arraybackslash}p{(\columnwidth - 2\tabcolsep) * \real{0.5000}}
  >{\raggedright\arraybackslash}p{(\columnwidth - 2\tabcolsep) * \real{0.5000}}@{}}
\toprule\noalign{}
\begin{minipage}[b]{\linewidth}\raggedright
Dates in 2026
\end{minipage} & \begin{minipage}[b]{\linewidth}\raggedright
Milestone and ownership
\end{minipage} \\
\midrule\noalign{}
\endhead
\bottomrule\noalign{}
\endlastfoot
Oct 9-14 & Confirm platform, trace the dependency loop, and run a
limited static pilot. Chenrui: tasks/checks; Lingxiao:
resources/profiling. \\
Oct 15/20; report Oct 22, 9 a.m. & Midterm presentation slot and report:
revised scope, available evidence, unresolved access/measurement issues.
Both. \\
Oct 23-Nov 8 & Implement shared admission hooks and both online
policies; calibrate ranges/windows. Lingxiao: policies; Chenrui:
load/replay. \\
Nov 9-22 & Held-out stationary and transition experiments, ablations,
uncertainty analysis. Both. \\
Nov 23-Dec 8 & Clean reproduction, artifact release, and presentation
preparation. Confirm final presentation slot with course staff. Both. \\
Dec 9-15 & Final report and limitations; report due Dec 15, 9 a.m.
Both. \\
\end{longtable}

Plan 8-10 hours/person/week. Reduce operating points if compute is
limited; retain live validation and report policy overhead or lack of
benefit.

\hypertarget{relation-to-prior-work}{%
\subsection{Relation to prior work}\label{relation-to-prior-work}}

The course readings {[}1-3{]} motivate profiling, application-level
headroom, and workload-sensitive management. GrandSLAm {[}4{]} motivates
stage-level latency management; \(\mu\)Tune {[}5{]} motivates load
adaptation. Shenango {[}6{]} reallocates CPU cores, while Splitwise
{[}7{]} separates inference phases; our controller does neither. Raj et
al.~{[}8{]} already study agent CPU-aware batching/scheduling, and
AgentSysBench {[}9{]} studies heterogeneous components and deployment
sensitivity. We will reuse compatible artifacts. Our scoped contribution
is a reproducible comparison of global versus stage-specific online
control across measured resource conditions, without claiming that prior
work overlooks this dependence.

\hypertarget{references}{%
\subsection{References}\label{references}}

{[}1{]} Kanev et al.~\textbf{Profiling a Warehouse-Scale Computer.} ISCA
2015.
\href{https://research.google/pubs/profiling-a-warehouse-scale-computer/}{Publication}.

{[}2{]} Sriraman and Dhanotia. \textbf{Accelerometer: Understanding
Acceleration Opportunities for Data Center Overheads at Hyperscale.}
ASPLOS 2020.
\href{https://github.com/akshithasriraman/Accelerometer}{Artifact}.

{[}3{]} Shahrad et al.~\textbf{Serverless in the Wild: Characterizing
and Optimizing the Serverless Workload at a Large Cloud Provider.}
USENIX ATC 2020.
\href{https://www.usenix.org/conference/atc20/presentation/shahrad}{Paper}.

{[}4{]} Kannan et al.~\textbf{GrandSLAm: Guaranteeing SLAs for Jobs in
Microservices Execution Frameworks.} EuroSys 2019.
\href{https://ramsrivatsa.github.io/papers/kannan_eurosys19.pdf}{Paper}.

{[}5{]} Sriraman and Wenisch. \textbf{\(\mu\)Tune: Auto-Tuned Threading
for OLDI Microservices.} OSDI 2018.
\href{https://www.usenix.org/conference/osdi18/presentation/sriraman}{Paper}.

{[}6{]} Ousterhout et al.~\textbf{Shenango: Achieving High CPU
Efficiency for Latency-sensitive Datacenter Workloads.} NSDI 2019.
\href{https://www.usenix.org/conference/nsdi19/presentation/ousterhout}{Paper}.

{[}7{]} Patel et al.~\textbf{Splitwise: Efficient Generative LLM
Inference Using Phase Splitting.} ISCA 2024.
\href{https://www.microsoft.com/en-us/research/publication/splitwise-efficient-generative-llm-inference-using-phase-splitting/}{Paper}.

{[}8{]} Raj et al.~\textbf{Towards Understanding, Analyzing, and
Optimizing Agentic AI Execution: A CPU-Centric Perspective.}
arXiv:2511.00739v3, 2026 revision.
\href{https://arxiv.org/abs/2511.00739v3}{Paper}.

{[}9{]} Chang et al.~\textbf{From LLM Inference to Agentic Workloads:
Characterization and Implications for Serving Systems.} AgentSysBench,
arXiv:2608.15127, 2026. \href{https://arxiv.org/abs/2608.15127}{Paper}.

\textbf{AI assistance disclosure:} ChatGPT Codex assisted with research,
design, writing, and revision; the prior draft disclosed Claude
critique/editing. No results are asserted. Authors must review claims
and disclose subsequent AI use.

\end{document}
